\documentclass[10pt,a4paper]{amsart}
\usepackage[margin=19mm]{geometry}
\usepackage{amsmath,amssymb,mathtools}
\usepackage[scr=rsfs,cal=euler]{mathalfa}
\usepackage{xfrac}
\usepackage[unicode,hidelinks,bookmarks=false]{hyperref}
\newcommand{\ie}{i.e.\ }
\newcommand{\cf}{cf.\ }
\newcommand{\resp}{respectively\ }
\newcommand{\Subsection}{\subsection}
\providecommand{\nobreakdash}{\nobreak\hbox{-}}
\setlength{\emergencystretch}{2em}
\multlinegap0pt
\usepackage{amssymb}
\usepackage{mathtools}
\newtheorem{proposition}{Proposition}
\newcommand{\C}{\mathbb{C}}
\newcommand{\QLS}{\mathrm{QLS}}
\newcommand{\card}{\operatorname{card}}
\newcommand{\vectset}{\mathcal{V}}
\title{Three Quantum Latin Squares of Order 6 with Cardinalities 13, 15, and 17}
\author{Zhipeng Xu}
\date{Source: arXiv:2605.15540v2 (CC BY 4.0)}
\begin{document}
\maketitle

\noindent\emph{Source and licence.} Three Quantum Latin Squares of Order 6 with Cardinalities 13, 15, and 17. Version arXiv:2605.15540v2 (CC BY 4.0), distributed by arXiv under the Creative Commons Attribution 4.0 International licence, http://creativecommons.org/licenses/by/4.0/, as recorded in the arXiv licence metadata for this exact version and shown on the abstract page https://arxiv.org/abs/2605.15540v2. Full licence notice: \texttt{licenses/arxiv-2605.15540-CC-BY-4.0.txt}.\par\smallskip

\noindent\textit{Editorial scope.} Counted unit: the article's proposition proposition (source0.tex lines 176--178). The article's complete original proof of it is delivered with it (source0.tex lines 180--266, one \texttt{proof} environment of 87 lines). No mathematics is written, repaired or re-derived: the statement and the proof are the article's own lines, in the article's own notation, and no retained line is substituted or re-flowed.  No further source-local context is needed: every cross-reference of every retained line resolves inside the two delivered spans.  Nothing was omitted from any retained span, so the package carries no omission record.  No retained line contains a citation, so the wrapper carries no bibliography and no bibliography entry.  No retained line contains a figure environment, an image-inclusion command or drawing code, so no image is delivered and the package declares no asset dependency.  Editorial formatting only, in addition: an AMS \texttt{amsart} preamble that restates the article's own declarations and macros used by the retained lines, namely the environments \texttt{proposition} on separate counters, together with the standard packages those lines need.  The retained environments print in this document as Proposition 1 in order of appearance, whereas the article prints the same items with the numbering of its own full document.  The complete native source of the article as distributed in the arXiv e-print for this exact version is bundled unchanged as \texttt{sources/200664/source0.tex}.
\begin{proposition}
The array $\Phi_{13}$ is a $\QLS(6)$ with cardinality $13$.
\end{proposition}

\begin{proof}
Define three orthonormal bases of $U$ by
\[
  B_0=\{a,b,c,d\},
  \qquad
  B_1=\{a,e,f,g\},
  \qquad
  B_2=\{b,e,h,r\}.
\]
The set $B_0$ is the computational basis of $U$. Since
\[
  e=\frac{c+d}{\sqrt2},
  \qquad
  v=\frac{c-d}{\sqrt2},
\]
the pair $\{e,v\}$ is an orthonormal basis of $\operatorname{span}\{c,d\}$.
The pair $\{f,g\}$ is obtained from $\{b,v\}$ by the Hadamard rotation
\[
  f=\frac{b+v}{\sqrt2},
  \qquad
  g=\frac{b-v}{\sqrt2}.
\]
Therefore $\{f,g\}$ is an orthonormal basis of $\operatorname{span}\{b,v\}$,
and it is orthogonal to both $a$ and $e$. Hence $B_1$ is an orthonormal basis of
$U$.

Similarly, the pair $\{h,r\}$ is obtained from $\{a,v\}$ by the Hadamard
rotation
\[
  h=\frac{a+v}{\sqrt2},
  \qquad
  r=\frac{a-v}{\sqrt2}.
\]
Thus $\{h,r\}$ is an orthonormal basis of $\operatorname{span}\{a,v\}$, and it
is orthogonal to both $b$ and $e$. Hence $B_2$ is also an orthonormal basis of
$U$.

Define two orthonormal bases of $V$ by
\[
  S_0=\{x,y\},
  \qquad
  S_1=\{z,w\}.
\]
The set $S_0$ is the computational basis of $V$, and $S_1$ is obtained from
$S_0$ by a two-dimensional Hadamard rotation. Since $U\perp V$, every union
$B_i\cup S_j$ is an orthonormal basis of $\C^6=U\oplus V$.

The six rows of $\Phi_{13}$ are of the following types:
\[
\begin{array}{c|c}
\text{row} & \text{orthonormal basis} \\\hline
R_1 & B_0\cup S_0\\
R_2 & B_0\cup S_1\\
R_3 & B_1\cup S_0\\
R_4 & B_1\cup S_1\\
R_5 & B_2\cup S_0\\
R_6 & B_2\cup S_1
\end{array}
\]
and the six columns are of the following types:
\[
\begin{array}{c|c}
\text{column} & \text{orthonormal basis} \\\hline
C_1 & B_0\cup S_0\\
C_2 & B_1\cup S_1\\
C_3 & B_1\cup S_0\\
C_4 & B_2\cup S_0\\
C_5 & B_2\cup S_1\\
C_6 & B_0\cup S_1
\end{array}
\]
Thus every row and every column is an orthonormal basis of $\C^6$, so
$\Phi_{13}$ is a $\QLS(6)$.

The vectors appearing in $\Phi_{13}$ are exactly
\[
  \vectset(\Phi_{13})=\{a,b,c,d,e,f,g,h,r,x,y,z,w\}.
\]
This set has $13$ elements. No two of them differ only by a global phase:
vectors from $U$ and vectors from $V$ have orthogonal supports; the four vectors
$x,y,z,w$ in $V$ are pairwise not scalar multiples; and the nine vectors
$a,b,c,d,e,f,g,h,r$ in $U$ have different supports or different relative signs
on a common support. Therefore
\[
  \card(\Phi_{13})=13.
\]
\end{proof}

\end{document}
